\documentclass[xcolor={usenames,svgnames}]{beamer}

\usepackage[francais]{babel}
\usepackage{rtxslides}
\usepackage{listings}
\usepackage{tikz}
\usetikzlibrary{shapes,fit}

\title{\rtx\ -- RMLL 2011}
\date{8 Juillet 2011}
\author{\rtx \\ \texttt{www.rathaxes.org}}

\definecolor{lightred}{RGB}{147,36,33}
\tikzset{componentarrow/.style={->, >=stealth, color=rathaxesred, ultra thick}}

\newcommand{\cemph}[1]{{\itshape{\textcolor{rathaxesred}{#1}}}}

\newcommand{\tred}[1]{\textcolor{rathaxesred}{#1}}

\tikzset{warrow/.style={->, >=stealth, color=white, ultra thick}}

\tikzset{graybox/.style={draw,rectangle,rounded corners=3pt,very thick,densely dashed,color=gray!75,text=white}}
\tikzset{redbox/.style={draw,rectangle,rounded corners=5pt,ultra thick,color=rathaxesred,text=white}}
\tikzset{redcontainer/.style={draw,rectangle,rounded corners=5pt,ultra thick,color=rathaxesred,text=white,minimum height=3.5cm,minimum width=2.5cm}}

\lstset{ %
numbers=left,                   % where to put the line-numbers
numberstyle=\footnotesize,      % the size of the fonts that are used for the line-numbers
stepnumber=0,                   % the step between two line-numbers. If it's 1, each line
                                % will be numbered
basicstyle=\footnotesize\ttfamily,
numbersep=5pt,                  % how far the line-numbers are from the code
backgroundcolor=\color{black},  % choose the background color. You must add \usepackage{color}
showspaces=false,               % show spaces adding particular underscores
showstringspaces=false,         % underline spaces within strings
showtabs=false,                 % show tabs within strings adding particular underscores
tabsize=2,                      % sets default tabsize to 2 spaces
captionpos=b,                   % sets the caption-position to bottom
breaklines=true,                % sets automatic line breaking
breakatwhitespace=false,        % sets if automatic breaks should only happen at whitespace
escapeinside={\%*}{*)},         % if you want to add a comment within your code
morekeywords={*,...}            % if you want to add more keywords to the set
}

\begin{document}

\begin{frame}
\titlepage
\end{frame}

\begin{frame}{Historique du projet}
\begin{center}
\Huge{TODO} % TODO
\end{center}
\end{frame}

\begin{frame}{Avancement du projet}
\begin{center}
\begin{tikzpicture}
\node (P) at (0,2.5) {\LARGE{pilotes}};
\node (C) at (0,-2.5) {\LARGE{\cemph{compilateur}}};
\draw (-5,0) node[above left] {\LARGE{Itérations :}};
\draw[warrow] (-1.25,2.165) arc(120:240:2.5);
\draw[warrow] (1.25,-2.165) arc(-60:60:2.5);
\end{tikzpicture}
\end{center}
\end{frame}

\begin{frame}{Pilotes de périphériques}
\begin{center}
\only<1>{\LARGE{Souvent \cemph{\Huge{instables}} et \cemph{\Huge{bancals}}}}
\only<2>{%
\large{\itshape{\rmfamily{«~\tred{30 \%} of Windows Vista \tred{crashes} caused by \tred{Nvidia drivers}~»}}}

\vspace{1cm}
\raggedleft\large{2.6.39: \itshape{\rmfamily{«~les pilotes (\tred{65 \% des patches})~»}} (patrick\_g)}

\vspace{1cm}
\large{2.6.35: \itshape{\rmfamily{«~les \tred{deux tiers des changements} dans les pilotes~»}}}
}
\end{center}
\end{frame}

\begin{frame}{Développement de pilotes de périphériques}
\begin{center}
\only<1>{\Large{Une tâche \cemph{\LARGE{difficile}}, \cemph{\LARGE{longue}} et \cemph{\LARGE{fastidieuse}}}}
\only<2>{\Large{%
\begin{itemize}
\item Compétences requises ;
\item Portabilité ;
\item Maintenance.
\end{itemize}
}}
\end{center}
\end{frame}

\begin{frame}{Objectifs du DSL \rtx}
\begin{center}
\only<1>{\Large{\cemph{\LARGE{Simplifier}} les cycles de développement}}
\only<2>{\Large{\begin{itemize}
\item Abstraction de la plateforme ;
\item Vérifications statiques ;
\item Séparation des compétences.
\end{itemize}}}
\end{center}
\end{frame}

\begin{frame}{Génération d'un pilote}
\only<1>{\begin{tikzpicture}[overlay]
\node (RTX) at (2,2.5) {\includegraphics[height=1.5cm]{icons/rtx}};
\node[redbox] (COMPILER) at (2,-0.5) {Compilateur};
\node (C) at (2,-3.5) {\includegraphics[height=1.5cm]{icons/c}};

\draw[warrow] (RTX)--(COMPILER);
\draw[warrow] (COMPILER)--(C);

\draw (RTX.north east) node[below right]{%
\begin{minipage}{9cm}
\large{Description du pilote :}
\small{
\begin{itemize}
\item Description du matériel ;
\item Algorithmes ;
\item Paramètres spécifiques à chaque système.
\end{itemize}}
\end{minipage}};

\draw (C.north east) node [below right]{\large{Module noyau}};
\end{tikzpicture}}

\only<2>{\begin{tikzpicture}[overlay]
\node (RTX) at (2,2.5) {\includegraphics[height=1.5cm]{icons/rtx}};
\node[redbox] (COMPILER) at (2,-0.5) {Compilateur};
\node (C) at (2,-3.5) {\includegraphics[height=1.5cm]{icons/c}};

\draw[warrow] (RTX)--(COMPILER);
\draw[warrow] (COMPILER)--(C);

\node[below right] (H1) at (RTX.north east) {%
\begin{minipage}{7.5cm}
\large{Description du pilote :} \\
\small{\texttt{
device Pipe use LKM, CharDev, Algorithms \{
\hspace*{1ex} open(Context ctx) \{ \\
\hspace*{2ex} log(“Open called on device Pipe”); \\
\hspace*{1ex} \} \\
\}
}}
\end{minipage}};
\draw (C.north east) node [below right]{\large{Module noyau}};
\end{tikzpicture}}

\only<3>{\begin{tikzpicture}[overlay]
\node[opacity=0.3] (RTX) at (2,2.5) {\includegraphics[height=1.5cm]{icons/rtx}};
\node[redcontainer] (COMPILER) at (2,-0.5) {};
\node[opacity=0.3] (C) at (2,-3.5) {\includegraphics[height=1.5cm]{icons/c}};

\draw[warrow,opacity=0.3] (RTX)--(COMPILER);
\draw[warrow,opacity=0.3] (COMPILER)--(C);

\draw (RTX.north east) node[below right, opacity=0.3]{\large{Description du pilote}};
\draw (COMPILER.south east) node[above right]{\Large{Compilateur}};
\draw (C.north east) node [below right, opacity=0.3]{\large{Module noyau}};
\end{tikzpicture}}

\only<4>{\begin{tikzpicture}[overlay]
\node[opacity=0.3] (RTX) at (2,2.5) {\includegraphics[height=1.5cm]{icons/rtx}};
\draw (RTX.north east) node[below right, opacity=0.3]{\large{Description du pilote}};

\node[opacity=0.3] (C) at (2,-3.5) {\includegraphics[height=1.5cm]{icons/c}};

\node[graybox] (AST1) at (2,0.5){\texttt{AST}};
\draw[warrow] (RTX)--(AST1);

%\node[graybox] (AST2) at (2,-1.5) {\texttt{AST}};
%\draw[warrow] (AST2)--(C);

\node[redcontainer] (COMPILER) at (2,-0.5) {};
\draw (COMPILER.south east) node[above right]{\Large{Compilateur}};

%\node[graybox] (RTXLINK) at (COMPILER.center) {\texttt{RTX\_Link}};
%\draw[warrow] (AST1)--(RTXLINK);
%\draw[warrow] (RTXLINK)--(AST2);

\draw[warrow,opacity=0.3] (COMPILER)--(C);
\draw (C.north east) node [below right, opacity=0.3]{\large{Module noyau}};
\end{tikzpicture}}

\only<5>{\begin{tikzpicture}[overlay]
\node[opacity=0.3] (RTX) at (2,2.5) {\includegraphics[height=1.5cm]{icons/rtx}};
\draw (RTX.north east) node[below right, opacity=0.3]{\large{Description du pilote}};

\node[opacity=0.3] (C) at (2,-3.5) {\includegraphics[height=1.5cm]{icons/c}};

\node[graybox] (AST1) at (2,0.5){\texttt{AST}};
\draw[warrow] (RTX)--(AST1);

%\node[graybox] (AST2) at (2,-1.5) {\texttt{AST}};
%\draw[warrow] (AST2)--(C);

\node[redcontainer] (COMPILER) at (2,-0.5) {};
\draw (COMPILER.south east) node[above right]{\Large{Compilateur}};

\node[graybox] (RTXLINK) at (COMPILER.center) {\texttt{RTX\_Link}};
\draw[warrow] (AST1)--(RTXLINK);
%\draw[warrow] (RTXLINK)--(AST2);

\draw[warrow,opacity=0.3] (COMPILER)--(C);
\draw (C.north east) node [below right, opacity=0.3]{\large{Module noyau}};
\end{tikzpicture}}

\only<6>{\begin{tikzpicture}[overlay]
\node[opacity=0.3] (RTX) at (2,2.5) {\includegraphics[height=1.5cm]{icons/rtx}};
\draw (RTX.north east) node[below right, opacity=0.3]{\large{Description du pilote}};

\node[opacity=0.3] (C) at (2,-3.5) {\includegraphics[height=1.5cm]{icons/c}};

\node[graybox] (AST1) at (2,0.5){\texttt{AST}};
\draw[warrow] (RTX)--(AST1);

%\node[graybox] (AST2) at (2,-1.5) {\texttt{AST}};
%\draw[warrow] (AST2)--(C);

\node[redcontainer] (COMPILER) at (2,-0.5) {};
\draw (COMPILER.south east) node[above right]{\Large{Compilateur}};

\node[graybox] (RTXLINK) at (COMPILER.center) {\texttt{RTX\_Link}};
\draw[warrow] (AST1)--(RTXLINK);
%\draw[warrow] (RTXLINK)--(AST2);

\draw[warrow,opacity=0.3] (COMPILER)--(C);

\node (RTI) at (4.5,0.85) {\includegraphics[height=1.3cm]{icons/rti}};
\draw (RTI.east) node[above right]{\Large{Interfaces}};
\draw[warrow] (RTI.west)--(RTXLINK.north);
\draw (C.north east) node [below right, opacity=0.3]{\large{Module noyau}};
\end{tikzpicture}}

\only<7>{\begin{tikzpicture}[overlay]
\node[opacity=0.3] (RTX) at (2,2.5) {\includegraphics[height=1.5cm]{icons/rtx}};
\draw (RTX.north east) node[below right, opacity=0.3]{\large{Description du pilote}};

\node[opacity=0.3] (C) at (2,-3.5) {\includegraphics[height=1.5cm]{icons/c}};

\node[graybox] (AST1) at (2,0.5){\texttt{AST}};
\draw[warrow] (RTX)--(AST1);

%\node[graybox] (AST2) at (2,-1.5) {\texttt{AST}};
%\draw[warrow] (AST2)--(C);

\node[redcontainer] (COMPILER) at (2,-0.5) {};
\draw (COMPILER.south east) node[above right]{\Large{Compilateur}};

\node[graybox] (RTXLINK) at (COMPILER.center) {\texttt{RTX\_Link}};
\draw[warrow] (AST1)--(RTXLINK);
%\draw[warrow] (RTXLINK)--(AST2);

\draw[warrow,opacity=0.3] (COMPILER)--(C);

\node (RTI) at (4.5,0.85) {\includegraphics[height=1.3cm]{icons/rti}};
\draw (RTI.east) node[above right]{\Large{Interfaces}};
\draw[warrow] (RTI.west)--(RTXLINK.north);

\node (BLT) at (4.5,-0.5) {\includegraphics[height=1.3cm]{icons/blt}};
\draw (BLT.east) node[above right]{\Large{Patrons de code}};
\draw[warrow] (BLT.west)--(RTXLINK);
\draw (C.north east) node [below right, opacity=0.3]{\large{Module noyau}};
\end{tikzpicture}}

\only<8>{\begin{tikzpicture}[overlay]
\node[opacity=0.3] (RTX) at (2,2.5) {\includegraphics[height=1.5cm]{icons/rtx}};
\draw (RTX.north east) node[below right, opacity=0.3]{\large{Description du pilote}};

\node[opacity=0.3] (C) at (2,-3.5) {\includegraphics[height=1.5cm]{icons/c}};

\node[graybox] (AST1) at (2,0.5){\texttt{AST}};
\draw[warrow] (RTX)--(AST1);

\node[graybox] (AST2) at (2,-1.5) {\texttt{AST}};
\draw[warrow] (AST2)--(C);

\node[redcontainer] (COMPILER) at (2,-0.5) {};
\draw (COMPILER.south east) node[above right]{\Large{Compilateur}};

\node[graybox] (RTXLINK) at (COMPILER.center) {\texttt{RTX\_Link}};
\draw[warrow] (AST1)--(RTXLINK);
\draw[warrow] (RTXLINK)--(AST2);

\node (RTI) at (4.5,0.85) {\includegraphics[height=1.3cm]{icons/rti}};
\draw (RTI.east) node[above right]{\Large{Interfaces}};
\draw[warrow] (RTI.west)--(RTXLINK.north);

\node (BLT) at (4.5,-0.5) {\includegraphics[height=1.3cm]{icons/blt}};
\draw (BLT.east) node[above right]{\Large{Patrons de code}};
\draw[warrow] (BLT.west)--(RTXLINK);
\draw (C.north east) node [below right, opacity=0.3]{\large{Module noyau}};
\end{tikzpicture}}
\end{frame}

\begin{frame}{Interfaces et patrons de code}
\begin{tikzpicture}[overlay]
\node (RTI) at (1.2,1.6) {\includegraphics[height=1.5cm]{icons/rti}};

\draw (RTI.north east) node[below right]{%
\begin{minipage}{10cm}
Décrivent un \emph{«~sous système~»} avec :
\begin{itemize}
\item Des types ;
\item Des séquences ;
\item Des variables de configuration.
\end{itemize}
\end{minipage}};
\end{tikzpicture}
\end{frame}

\begin{frame}[fragile]{Interfaces et patrons de code}
\begin{tikzpicture}[overlay]
\node (RTI) at (1.2,1.6) {\includegraphics[height=1.5cm]{icons/rti}};

\draw (RTI.east) node[right]{%
\begin{minipage}{10cm}
\begin{lstlisting}
interface CharDev : LKM {
  provided type        Context;
  required sequence    open(Context) {
     use pointcut      LKM::GLOBAL_CODE_DEFINITION;
     use pointcut      LKM::INIT_LKM_FPTRS;
     provided pointcut ALGO;
  }
}
\end{lstlisting}
\end{minipage}};
\end{tikzpicture}
\end{frame}

\begin{frame}[fragile]{Interfaces et patrons de code}
\begin{tikzpicture}[overlay]
\node (RTI) at (1.2,1.6) {\includegraphics[height=1.5cm]{icons/rti}};

\draw (RTI.east) node[right]{%
\begin{minipage}{10cm}
\begin{lstlisting}
interface CharDev : LKM {
  provided type        Context;
  required sequence    open(Context) {
     use pointcut      LKM::GLOBAL_CODE_DEFINITION;
     use pointcut      LKM::INIT_LKM_FPTRS;
     provided pointcut ALGO;
  }
}
\end{lstlisting}
\end{minipage}};

\node (BLT) at (1.2,-2.2) {\includegraphics[height=1.5cm]{icons/blt}};

\draw (BLT.north east) node[below right]{%
\begin{minipage}{10cm}
Implémentations des interfaces pour chaque OS :
\begin{itemize}
\item Templates de types ;
\item Templates de séquences.
\end{itemize}
\end{minipage}};
\end{tikzpicture}
\end{frame}

\begin{frame}[fragile]{Interfaces et patrons de code}
\begin{tikzpicture}[overlay]
\node (RTI) at (1.2,1.6) {\includegraphics[height=1.5cm]{icons/rti}};

\draw (RTI.east) node[right]{%
\begin{minipage}{10cm}
\begin{lstlisting}
interface CharDev : LKM {
  provided type        Context;
  required sequence    open(Context) {
     use pointcut      LKM::GLOBAL_CODE_DEFINITION;
     use pointcut      LKM::INIT_LKM_FPTRS;
     provided pointcut ALGO;
  }
}
\end{lstlisting}
\end{minipage}};

\node (BLT) at (1.2,-2.2) {\includegraphics[height=1.5cm]{icons/blt}};

\draw (BLT.east) node[right]{%
\begin{minipage}{10cm}
\begin{lstlisting}
with CharDev
values OS = Linux, version >= 2.6.27 {
  template sequence CharDev::open(Context ctx) {
     chunk LKM::GLOBAL_CODE_DEFINITION {
         /* C instrumente */
     }
  }
}
\end{lstlisting}
\end{minipage}};
\end{tikzpicture}
\end{frame}

\begin{frame}[fragile]{Tissage des ASTs}
\begin{tikzpicture}[overlay]
\node (RTX) at (1.2,2.4) {\includegraphics[height=1.5cm]{icons/rtx}};
\draw (RTX.east) node[right]{%
\begin{minipage}{115mm}
\begin{lstlisting}
device Pipe use LKM, CharDev, Algorithms {
  open(Context ctx) {
     log("Open called on device Pipe");
  }
}
\end{lstlisting}
\end{minipage}};

\node (RTI) at (1.2,-0.1) {\includegraphics[height=1.5cm]{icons/rti}};

\draw (RTI.east) node[right]{%
\begin{minipage}{115mm}
\begin{lstlisting}
interface CharDev : LKM {
  provided type        Context;
  required sequence    open(Context) {
     use pointcut      LKM::GLOBAL_CODE_DEFINITION;
     use pointcut      LKM::INIT_LKM_FPTRS;
     provided pointcut ALGO;
  }
}
\end{lstlisting}
\end{minipage}};

\node (BLT) at (1.2,-3.2) {\includegraphics[height=1.5cm]{icons/blt}};

\draw (BLT.east) node[right]{%
\begin{minipage}{115mm}
\begin{lstlisting}
with CharDev
values OS = Linux, version >= 2.6.27 {
  template sequence CharDev::open(Context ctx) {
     chunk LKM::GLOBAL_CODE_DEFINITION {
         ${pointcut ALGO}
     }
  }
}
\end{lstlisting}
\end{minipage}};
\end{tikzpicture}
\end{frame}

\begin{frame}[fragile]{Tissage des ASTs}
\begin{tikzpicture}[overlay]
\node (RTX) at (1.2,2.4) {\includegraphics[height=1.5cm]{icons/rtx}};
\draw (RTX.east) node[right]{%
\begin{minipage}{115mm}
\begin{lstlisting}
device Pipe use LKM, CharDev, Algorithms {
  open(Context ctx) {
     log("Open called on device Pipe");
  }
}
\end{lstlisting}
\end{minipage}};

\node (RTI) at (1.2,-0.1) {\includegraphics[height=1.5cm]{icons/rti}};

\draw (RTI.east) node[right]{%
\begin{minipage}{115mm}
\begin{lstlisting}
interface CharDev : LKM {
  provided type        Context;
  required sequence    open(Context) {
     use pointcut      LKM::GLOBAL_CODE_DEFINITION;
     use pointcut      LKM::INIT_LKM_FPTRS;
     provided pointcut ALGO;
  }
}
\end{lstlisting}
\end{minipage}};

\node (BLT) at (1.2,-3.2) {\includegraphics[height=1.5cm]{icons/blt}};

\draw (BLT.east) node[right]{%
\begin{minipage}{115mm}
\begin{lstlisting}
with CharDev
values OS = Linux, version >= 2.6.27 {
  template sequence CharDev::open(Context ctx) {
     chunk LKM::GLOBAL_CODE_DEFINITION {
         ${pointcut ALGO}
     }
  }
}
\end{lstlisting}
\end{minipage}};

\draw [ultra thick, rounded corners=3pt, color=LimeGreen] (2.5,0.27) rectangle (9.15,0.67);
\draw [ultra thick, color=LimeGreen, ->, >=stealth] (7,0.67) -- (3.1,2.65);
\draw [ultra thick, color=LimeGreen, ->, >=stealth] (7,0.27) -- (8,-2.50);
\end{tikzpicture}
\end{frame}

\begin{frame}[fragile]{Tissage des ASTs}
\begin{tikzpicture}[overlay]
\node (RTX) at (1.2,2.4) {\includegraphics[height=1.5cm]{icons/rtx}};
\draw (RTX.east) node[right]{%
\begin{minipage}{115mm}
\begin{lstlisting}
device Pipe use LKM, CharDev, Algorithms {
  open(Context ctx) {
     log("Open called on device Pipe");
  }
}
\end{lstlisting}
\end{minipage}};

\node (RTI) at (1.2,-0.1) {\includegraphics[height=1.5cm]{icons/rti}};

\draw (RTI.east) node[right]{%
\begin{minipage}{115mm}
\begin{lstlisting}
interface CharDev : LKM {
  provided type        Context;
  required sequence    open(Context) {
     use pointcut      LKM::GLOBAL_CODE_DEFINITION;
     use pointcut      LKM::INIT_LKM_FPTRS;
     provided pointcut ALGO;
  }
}
\end{lstlisting}
\end{minipage}};

\node (BLT) at (1.2,-3.2) {\includegraphics[height=1.5cm]{icons/blt}};

\draw (BLT.east) node[right]{%
\begin{minipage}{115mm}
\begin{lstlisting}
with CharDev
values OS = Linux, version >= 2.6.27 {
  template sequence CharDev::open(Context ctx) {
     chunk LKM::GLOBAL_CODE_DEFINITION {
         ${pointcut ALGO}
     }
  }
}
\end{lstlisting}
\end{minipage}};

\draw [ultra thick, rounded corners=3pt, color=LimeGreen] (2.5,0.27) rectangle (9.15,0.67);
\draw [ultra thick, color=LimeGreen, ->, >=stealth] (7,0.67) -- (3.1,2.65);
\draw [ultra thick, color=LimeGreen, ->, >=stealth] (7,0.27) -- (8,-2.50);

\draw [ultra thick, rounded corners=3pt, color=MediumOrchid] (3.1,-0.47) rectangle (7.6,-0.87);
\draw [ultra thick, color=MediumOrchid, ->, >=stealth] (6,-0.47) -- (5,2.25);
\draw [ultra thick, color=MediumOrchid, ->, >=stealth] (6,-0.87) -- (5,-3.27);
\end{tikzpicture}
\end{frame}

\begin{frame}{Tissage des ASTs}
\begin{tikzpicture}[overlay]
\node (RTX) at (1.2,2.4) {\includegraphics[height=1.5cm]{icons/rtx}};
\draw (RTX.east) node[right]{\LARGE{Arbre d'\emph{« advices »}}};

\node (RTI) at (1.2,-0.1) {\includegraphics[height=1.5cm]{icons/rti}};

\draw (RTI.east) node[right]{\LARGE{\emph{« Abstraction »} du backend}};

\node (BLT) at (1.2,-3.2) {\includegraphics[height=1.5cm]{icons/blt}};

\draw (BLT.east) node[right]{\LARGE{Arbre de \emph{« pointcuts »}}};
\end{tikzpicture}
\end{frame}

\begin{frame}{Questions}
\begin{center}
\Huge{Merci}

\end{center}

%\hspace{1em}\rule{5cm}{0.2mm}
\vspace{2em}
\begin{itemize}
\item \Large{\texttt{http://www.rathaxes.org/}}
\item \Large{\texttt{\#rathaxes} sur \texttt{irc.freenode.net}}
\end{itemize}
\end{frame}

\end{document}
